
%%%%%%%%%%%%% SI!!!! %%%%%%%%%%%%%%%
\ifastroph
   \setcounter{figure}{0}
   \renewcommand{\thefigure}{S\arabic{figure}}
\else
  \ifsuppl
     \setcounter{figure}{-4}
  \fi
\fi


\fignature{3Dplot.jpg}{0.98}
{3D-plot of the light path. North is toward
the positive $y$-axis (up), east is toward the negative $x$-axis
(left), and the positive $z$-axis points toward the observer with the
origin at $\eta$~Car.  The red and brown circles indicate $\eta$~Car
and scattering dust, respectively.  The parabolic relation between the
spatial parameters of the scattering dust and the time since outburst
is described by the well-known light echo equation\cite{Couderc39}:
Assuming a time since outburst of 169 years, and a distance of 7660
light-years\cite{Smith06}, we find that the scattering dust is at a
position ($x$,$y$,$z$)=(14,-78,-66) in light-years.  
%The $\eta$~Car to
%light-echo cone apex distance, in light years, is half the interval of
%time since the $\eta$~Car outburst. 
The black lines show the path of
the light scattering from the light-echo-producing dust
concentrations.  The top-right panel shows an \textit{HST\/} image
(credit: Nathan Smith / Jon Morse / NASA) of $\eta$~Car, which is
surrounded by expanding lobes of gas denoted as the Homunculus Nebula
visible to the NW and SE. These lobes were created by the 1838-1858
Great Eruption. We indicate the lobes in the 3D plots with the two
cones (Note that the size of the cones is not to scale).
}
{fig:3Dplot}

\fignature{lines_3.pdf}{0.57}
{  Spectral
  type comparison. The stellar comparison spectra are convolved with a
  Gaussian of FWHM 7~\AA\ so that their resolution matches those 
  of \specMlo\ and \specDP. Similarly, we convolve
  the spectrum \specMhi, which has an original resolution of 4\AA,
  with a Gaussian of FWHM 5.7~\AA. We then divide the spectra by
  the low-order continuum which we determine by convolving the spectra
  with a Gaussian of FWHM 200~\AA. The upper left panel
  compares two of the observed light echo
  spectra to a selection of UVES supergiants\cite{Bagnulo03} in the 5060-5500~\AA\
  wavelength range (We do not show spectrum \specMhi\ since its S/N
  ratio in this blue wavelength range is too low).  The light echo
  spectra correlate very well with late-F and G-type stars, in
  particular the Mg~b lines, and the \ion{Ca}{I}, \ion{Fe}{I},
  \ion{Ti}{I}, \ion{Cr}{I} blend at 5270~\AA. The upper right panel
  shows the cross-correlation between the light echo and the UVES
  spectra in the wavelength range of 5050-6500~\AA\  
  as calculated by
  the IRAF routine {\tt xcsao}.
%
%, and the \ion{Fe}{I},
%  \ion{Ba}{I}, \ion{Ca}{I}, \ion{Mn}{I}, \ion{Co}{I},
%  \ion{Ti}{I}~and~II, \ion{Ni}{I} blend at 6497~\AA (\specMlo\ has a
%  chip gap at the 6497~\AA\ blend). The \ion{Na}{I}~D resonance
%  doublet near 5890~\AA\ is prominent in both the echo and comparison
%  spectra. Unfortunately, it suffers from contamination by the
%  \ion{He}{I} $\lambda$5876 emission from the Carina Nebula at the
%  blueward edge; extremely strong, multiple-component interstellar
%  \ion{Na}{I}~D absorption within the Carina Nebula\cite{Walborn07};
%  and residuals from night sky emission at the redward edge. 
  Since the UVES spectra have a gap around the \ion{Ca}{II}
  IR triplet, we use the \ion{Ca}{II} triplet spectral
  library\cite{Cenarro01} in that wavelength range, as shown in the
  bottom left panel.
%In order to guard against biasing the continuum by the strong \ion{Ca}{II}
%IR triplet lines, we disregard the wavelength ranges
%8484-8505~\AA, 8520-8546~\AA, and 8643-8670~\AA\ when we compute the
%continuum.  
%  The bottom middle panel shows the correlation $r$
%  between the light echo and the UVES spectra in the wavelength range
%  of 5050-6500~\AA\ as calculated by the IRAF routine {\tt xcsao}. 
%We cross-correlate each of our $\eta$~Car spectra with the supergiant
%spectra library in Fourier space using the {\tt xcsao}\cite{Kurtz92}
%IRAF task in order to determine which stellar type most strongly
%correlates. Velocity shifts and dispersions are optimized for each
%spectrum. An initial velocity guess of
%$v_\mathrm{guess}=-250~\mathrm{km~s}^{-1}$ is provided, based on
%visual inspection of the \ion{Ca}{II} triplet, and a velocity range of
%$-350~\mathrm{km~s}^{-1}<v<-150~\mathrm{km~s}^{-1}$ is allowed for the
%fit. Spectra are processed separately in the range $5050$~\AA$\ <
%\lambda <6500$~\AA\ and in the range $8200$~\AA$\ < \lambda
%<8720$~\AA, giving consistent results within the errorbars. The
%\ion{Na}{I} D line region, $5860$~\AA$\ < \lambda < 5920$~\AA, and
%$5890$~\AA$\ < \lambda <5866$~\AA, is removed, as these lines are
%affected by ISM absorption, which contaminates the result. Together
%with the velocity, a correlation parameter $r$ is determined.
  The bottom right panel shows the blue-shifted velocity determined by {\tt
  xcsao} from the \ion{Ca}{II} triplet with respect to the effective temperature\cite{Humphreys84}
  of the supergiant template spectra. 
  %For each supergiant type, we
  %use the relation from\cite{Humphreys84} to determine the
  %temperature.
}
{fig:speclines}

\fignature{probdensity2.pdf}{0.63}
{Effective temperature of the light echo spectrum.  The left panel
shows the cross-correlation between the light echo and the UVES
spectra as calculated by the IRAF routine {\tt xcsao} for \specMlo\
(red squares) and \specDP\ (blue diamonds).  The lines indicate $r(T)$
smoothed with a Gaussian of width $\sigma=300$~K, peaking at 5210~K
and 4950~K for \specMlo\ (red) and \specDP\ (blue), respectively. The
right panel shows the PDF of the best-correlating supergiant
temperature, determined by bootstrap resampling the $r$ distribution
$10^5$~times. The 95\% temperature confidence intervals are
4850-5550~K and 4450-5400~K for \specMlo\ and \specDP, respectively.}
{fig:pdf}

